\subsection{压缩的条件尾平方函数与 Schatten 刻画}

沿用第 \ref{sec:martingale} 节的壳过滤，记
\[
 x_n=\mathbb E_nX,\qquad
 \Delta_1X=\mathbb E_1X,\qquad
 \Delta_kX=x_k-x_{k-1}\quad(k\geq2).
\]
这里 $\Delta_k$ 是 PSZ 的鞅差分算子 $d_k$；使用不同字母以区别于重数
$d_n=\rank E_n$。由 \eqref{eq:shell-difference} 的无限角版本，
\begin{equation}\label{sch:eq:difference}
 \Delta_kX=P_{k-1}XE_k+E_kXP_{k-1},\qquad k\geq2.
\end{equation}
这将 PSZ Section 2.3 的行、列平方函数用于壳鞅的未来差分。
定义压缩的条件尾平方函数
\begin{align}
 s_{n,c}(X)
 &=\left[P_n\mathbb E_n\bigl(|X-x_n|^2\bigr)P_n\right]^{1/2},
       \label{sch:eq:tail-column}\\
 s_{n,r}(X)
 &=\left[P_n\mathbb E_n\bigl(|X^*-x_n^*|^2\bigr)P_n\right]^{1/2}.
       \label{sch:eq:tail-row}
\end{align}
这两个正算子均支撑在有限秩投影 $P_n$ 上。
注意行型中 $|X^*-x_n^*|^2=(X-x_n)(X-x_n)^*$。
上式使用 $X-x_n$，即第 $n$ 步之后的鞅尾部；
它与 BMO 定义 \eqref{eq:martingale-bmoc} 中的 $X-x_{n-1}$ 有一个明确的指标区别。

\begin{theorem}[壳鞅尾平方函数的 Schatten 可和性]\label{sch:thm:main}
对任意 $0<p<\infty$（特别是 $1<p<\infty$），
\begin{equation}\label{sch:eq:criterion}
 \boxed{\quad
 \CC_X\in S_p\bigl(S_2(H)\bigr)
 \iff
 \sum_{n\geq1}d_n\left(
   \norm{s_{n,c}(X)}_{S_p(H)}^p+
   \norm{s_{n,r}(X)}_{S_p(H)}^p\right)<\infty.
 \quad}
\end{equation}
而且有精确等式
\begin{equation}\label{sch:eq:main}
 \norm{\CC_X}_{S_p(S_2(H))}^p
 =\sum_{n\geq1}d_n\left(
   \norm{s_{n,c}(X)}_{S_p(H)}^p+
   \norm{s_{n,r}(X)}_{S_p(H)}^p\right).
\end{equation}
两种平方函数还可以完全由未来鞅差分表示：
\begin{align}
 s_{n,c}(X)^2
 &=\sum_{k>n}(\Delta_kX\,P_n)^*(\Delta_kX\,P_n),
       \label{sch:eq:tail-diff-column}\\
 s_{n,r}(X)^2
 &=\sum_{k>n}(P_n\Delta_kX)(P_n\Delta_kX)^*.
       \label{sch:eq:tail-diff-row}
\end{align}
对固定 $n$，这些正算子级数在 $B(P_nH)$ 中按算子范数收敛。
等价地，保留原先的角块形式，仍有
\begin{equation}\label{sch:eq:corners}
 \norm{\CC_X}_{S_p(S_2(H))}^p
 =\sum_{n\geq1}d_n\norm{[X,P_n]}_{S_p(H)}^p
 =\sum_{n\geq1}d_n\left(
  \norm{(I-P_n)XP_n}_{S_p(H)}^p+
  \norm{P_nX(I-P_n)}_{S_p(H)}^p\right).
\end{equation}
对于 $0<p<1$，以上使用 Schatten 拟范数。
\end{theorem}

